%% ============================================================
%% Lambda-EG x JADES DR4 -- manuscrito para RMxAA
%% Clase rmxaa (RMxAA-journal/RMxAA-macros, v3)
%% Compilar:  pdflatex -> bibtex -> pdflatex -> pdflatex
%% ============================================================
\documentclass{rmxaa}

\usepackage[utf8]{inputenc}
\usepackage{amsmath}
\usepackage{paralist}
\usepackage[svgnames]{xcolor}
\hypersetup{colorlinks, linkcolor=MediumBlue, citecolor=black,
            urlcolor=MediumBlue}

%% --- Macros utiles ---
\newcommand{\gcrit}{g_{\rm crit}}
\newcommand{\Mbar}{M_{\rm bar}}
\newcommand{\Vobs}{V_{\rm obs}}
\newcommand{\Vpred}{V_{\rm pred}}
\newcommand{\kms}{\ensuremath{\mathrm{km\,s^{-1}}}}
\newcommand{\Ha}{\ensuremath{\mathrm{H\alpha}}}
\newcommand{\Msun}{\ensuremath{M_\odot}}
\newcommand{\LEG}{\mbox{$\Lambda$-EG}}

%% ============================================================
%% Encabezado
%% ============================================================
\title{A First Observational Test of the \LEG{} $\beta(\Phi)$ Flow
  against JADES DR4 NIRSpec Spectroscopy at $z=3$--$6$}

\author{J. P. Figueroa Torres\altaffilmark{1}}

\altaffiltext{1}{Independent Researcher, Guadalajara, Jalisco,
  M\'exico.}

\shortauthor{Figueroa Torres}
\shorttitle{A First Test of the \LEG{} $\beta(\Phi)$ Flow with JADES DR4}

\fulladdresses{
\item J. P. Figueroa Torres: Independent Researcher, Guadalajara,
  Jalisco, M\'exico (quimikcasera369@gmail.com).
}

\listofauthors{J. P. Figueroa Torres}
\indexauthor{Figueroa Torres, J. P.}

%% --- Abstract (ingles) ---
\abstract{Lambda Emergent Gravity (\LEG) predicts a redshift-dependent
  critical acceleration $\gcrit(z)=cH(z)/(4\pi^2\sqrt{3})$ and a
  baryonic Tully--Fisher relation (BTFR) that evolves as $V^4\propto
  \Mbar\,\gcrit(z)$. The framework further derives a ``$\beta$-function''
  $\beta(\Phi)$ governing the relaxation of the dimensionless BTFR
  deviation $\Phi$, with the quantitative prediction $\sigma_\Phi(0)<0.04$
  dex. We present the first quantitative observational test of this
  prediction using $59$ galaxies from the JWST/JADES Data Release~4 with
  NIRSpec $R\sim1000$ spectroscopy at $z=3$--$6$. We measure intrinsic
  \Ha{} velocity dispersions, estimate baryonic masses through
  \Ha$\to$SFR$\to$main-sequence calibrations, and evaluate $\Phi$ for the
  $24$ galaxies in the Taylor regime $|\Phi|<1$. The analytic shape of
  $\beta(\Phi)$ is recovered, but we find a systematic offset
  (median $\Vobs/\Vpred=0.77$, median $\Phi=-0.65$, $t$-test $p=0.014$)
  and $\sigma_\Phi(0)=0.10$ dex, exceeding the prediction by
  $\sim2.5\times$. A model-independent line-profile quality test
  (``SMITH-QA'') shows the offset is \emph{robust} to the exclusion of
  non-Gaussian profiles, ruling out turbulence/outflow contamination of
  $\sigma_{\Ha}$ as its origin. The bias corresponds to a $+0.45$~dex
  mass overestimate, which lies \emph{within} the $0.87$~dex inter-method
  scatter of high-$z$ stellar-mass calibrations. We conclude that the
  present test is \emph{calibration-limited} rather than theory-limited,
  in agreement with the framework's own forecast that a definitive test
  requires spatially resolved kinematics (ELT/HARMONI) and SED-fit
  masses.}

%% --- Resumen (espanol) ---
\resumen{La Gravedad Emergente Lambda (\LEG) predice una aceleraci\'on
  cr\'\i{}tica dependiente del corrimiento al rojo
  $\gcrit(z)=cH(z)/(4\pi^2\sqrt{3})$ y una relaci\'on de Tully--Fisher
  bari\'onica (BTFR) que evoluciona como $V^4\propto \Mbar\,\gcrit(z)$. El
  marco deriva adem\'as una ``funci\'on $\beta$'' $\beta(\Phi)$ que rige la
  relajaci\'on de la desviaci\'on adimensional $\Phi$ de la BTFR, con la
  predicci\'on cuantitativa $\sigma_\Phi(0)<0.04$~dex. Presentamos la
  primera prueba observacional cuantitativa de esta predicci\'on usando
  $59$ galaxias del Data Release~4 de JWST/JADES con espectroscop\'\i{}a
  NIRSpec $R\sim1000$ a $z=3$--$6$. Medimos dispersiones de velocidad
  intr\'\i{}nsecas de \Ha, estimamos masas bari\'onicas mediante
  calibraciones \Ha$\to$SFR$\to$secuencia principal, y evaluamos $\Phi$
  para las $24$ galaxias en el r\'egimen de Taylor $|\Phi|<1$. Se recupera
  la forma anal\'\i{}tica de $\beta(\Phi)$, pero hallamos un sesgo
  sistem\'atico (mediana $\Vobs/\Vpred=0.77$, mediana $\Phi=-0.65$,
  $p=0.014$) y $\sigma_\Phi(0)=0.10$~dex, que excede la predicci\'on por
  $\sim2.5\times$. Una prueba de calidad de perfil independiente del
  modelo (``SMITH-QA'') muestra que el sesgo es \emph{robusto} a la
  exclusi\'on de perfiles no gaussianos, descartando la contaminaci\'on
  por turbulencia/eyecciones como su origen. El sesgo corresponde a una
  sobrestimaci\'on de masa de $+0.45$~dex, que cae \emph{dentro} de la
  dispersi\'on entre m\'etodos de $0.87$~dex. Concluimos que la prueba
  est\'a \emph{limitada por la calibraci\'on} y no por la teor\'\i{}a, de
  acuerdo con la propia predicci\'on del marco.}

%% --- Keywords (lista estandar, orden alfabetico) ---
\addkeyword{cosmology: theory}
\addkeyword{dark matter}
\addkeyword{galaxies: evolution}
\addkeyword{galaxies: high-redshift}
\addkeyword{galaxies: kinematics and dynamics}
\addkeyword{gravitation}

\begin{document}
\maketitle

%% ============================================================
\section{Introduction}
\label{sec:intro}

The flat rotation curves of galaxies and the tight baryonic
Tully--Fisher relation (BTFR), $\Mbar\propto V^4$
\citep{McGaugh2000,LelliBTFR2016}, are usually attributed to dark-matter
halos, but they can also be read as evidence for a low-acceleration
modification of dynamics \citep{Milgrom1983} or for an emergent
gravitational response of the form proposed by \citet{Verlinde2017}.
A common feature of these scenarios is a \emph{critical acceleration}
scale $\gcrit\sim cH_0$ below which the effective dynamics depart from
the Newtonian regime.

Lambda Emergent Gravity \citep[\LEG;][]{FigueroaTorres2026} is a screened
scalar--tensor framework in which this scale arises from the spectral
geometry of the de~Sitter horizon, yielding the parameter-free value
\begin{equation}
  \label{eq:gcrit}
  \gcrit=\frac{cH_0}{4\pi^2\sqrt{3}}=9.58\times10^{-12}\
  \mathrm{m\,s^{-2}},
\end{equation}
with $H_0=67.4\ \kms\,{\rm Mpc^{-1}}$ \citep{Planck2020}. Because
$\gcrit$ inherits the Hubble rate, \LEG{} predicts a \emph{redshift
  evolution} of the BTFR through $\gcrit(z)=cH(z)/(4\pi^2\sqrt{3})$,
and Section~10 of \citet{FigueroaTorres2026} derives a renormalization
``$\beta$-function'' that governs how the dimensionless BTFR deviation
relaxes towards zero. This $\beta(\Phi)$ flow makes a sharp, falsifiable
prediction for the intrinsic scatter of the high-redshift BTFR
(\S\ref{sec:framework}).

The advent of JWST/NIRSpec spectroscopy of $z>3$ galaxies makes a direct
test of this prediction possible for the first time. In this Letter we
confront the \LEG{} $\beta(\Phi)$ flow with $59$ galaxies from the JADES
Data Release~4 (\S\ref{sec:data}), describe our spectral and mass
measurements (\S\ref{sec:methods}), present the test
(\S\ref{sec:results}), and ask whether the systematic offset we find is
astrophysical or instrumental by means of a model-independent
line-profile diagnostic (\S\ref{sec:smith}). We discuss the
implications in \S\ref{sec:discussion} and conclude in
\S\ref{sec:conclusions}. Throughout we adopt a flat $\Lambda$CDM
cosmology with $\Omega_m=0.315$, $\Omega_\Lambda=0.685$ and
$H_0=67.4\ \kms\,{\rm Mpc^{-1}}$ \citep{Planck2020}.

%% ============================================================
\section{The \LEG{} framework and the $\beta(\Phi)$ flow}
\label{sec:framework}

In \LEG{} the baryonic Tully--Fisher relation acquires an explicit
cosmological time dependence,
\begin{equation}
  \label{eq:btfr}
  V^4 = G\,\Mbar\,\gcrit(z), \qquad
  \gcrit(z)=\frac{c\,H(z)}{4\pi^2\sqrt{3}},
\end{equation}
with $H(z)=H_0\sqrt{\Omega_m(1+z)^3+\Omega_\Lambda}$. Equation
(\ref{eq:btfr}) implies that, at fixed baryonic mass, the rotation
velocity scales as
\begin{equation}
  \label{eq:Vevol}
  \frac{V(z)}{V(0)}=\left[\frac{H(z)}{H_0}\right]^{1/4},
\end{equation}
so that the BTFR zero-point brightens by $\approx0.16$ dex at $z=4$ and
$\approx0.20$ dex at $z=6$ relative to $z=0$.

It is convenient to define the dimensionless deviation of an individual
galaxy from the \LEG{} BTFR,
\begin{equation}
  \label{eq:Phi}
  \Phi \equiv \frac{\Vobs^4}{G\,\Mbar\,\gcrit(z)}-1
        = \left(\frac{\Vobs}{\Vpred}\right)^4-1,
\end{equation}
where $\Vpred=(G\,\Mbar\,\gcrit(z))^{1/4}$. Section~10 of
\citet{FigueroaTorres2026} shows that, under the screened $k$-essence
dynamics of the framework, $\Phi$ obeys a gradient flow whose exact
$\beta$-function is
\begin{equation}
  \label{eq:beta}
  \beta(\Phi)= -\,\frac{\Phi\,(2+\Phi)(1+\Phi)}{2+\Phi(2+\Phi)},
\end{equation}
valid for all $\Phi>-1$. Near the fixed point $\Phi=0$ this linearizes
to $\beta(\Phi)\simeq-\Phi-\tfrac12\Phi^2$, giving an exponential
relaxation $\Phi(\lambda)=\Phi_0\,e^{-\lambda}$ with flow time
$\lambda=\ln[H(z)/H_0]$. Consequently the dispersion of $\Phi$ across a
galaxy population is predicted to scale as
\begin{equation}
  \label{eq:sigmaPhi}
  \sigma_\Phi(z)=\sigma_\Phi(0)\,\frac{H(z)}{H_0},\qquad
  \sigma_\Phi(0)<0.04\ \mathrm{dex},
\end{equation}
the inequality being the quantitative target derived in \S10.8 of
\citet{FigueroaTorres2026}. Equations (\ref{eq:beta}) and
(\ref{eq:sigmaPhi}) are the two predictions we test below: the
\emph{shape} of the flow and the \emph{magnitude} of the residual
scatter.

%% ============================================================
\section{Data: the JADES DR4 sample}
\label{sec:data}

We use the JWST Advanced Deep Extragalactic Survey
\citep[JADES;][]{Bunker2024} Data Release~4 NIRSpec catalogue
\citep{CurtisLake2025,Scholtz2025}, specifically the combined external
table \texttt{Combined\_DR4\_external\_v1.2.1.fits}. From the parent
catalogue of $2{,}858$ sources with secure spectroscopic redshifts
(quality flag A) we select galaxies that

\begin{compactenum}[(i)]
\item lie at $3.0\le z\le 6.5$, where \Ha{} falls in the G235M/G395M
  gratings;
\item have a catalogue \Ha{} signal-to-noise ratio ${\rm SNR}>10$; and
\item have a well-defined rest-frame UV magnitude $-30<M_{\rm UV}<-10$.
\end{compactenum}

For these objects we retrieved the $R\sim1000$ one-dimensional extracted
spectra (\texttt{x1d}) from MAST. After the spectral quality cuts
described in \S\ref{sec:methods} we obtain a final working sample of
$59$ galaxies with intrinsically resolved \Ha{} dispersions.
Figure~\ref{fig:sample} shows their redshift distribution, which peaks
at $z\simeq3.5$ and extends to $z\simeq6$.

\begin{figure}[!t]
  \centering
  \includegraphics[width=\columnwidth]{fig1_sample}
  \caption{Spectroscopic redshift distribution of the $59$ JADES DR4
    galaxies in the working sample. \Ha{} is accessible with the
    NIRSpec medium-resolution gratings over the full $z=3$--$6$ range.}
  \label{fig:sample}
\end{figure}

%% ============================================================
\section{Methods}
\label{sec:methods}

\subsection{Intrinsic \Ha{} dispersions}
\label{sec:spec}

For each source we extract a $\pm0.05\,\mu$m window around the observed
\Ha{} wavelength $\lambda_{\rm obs}=6563\,(1+z)$\,\AA{} and fit a single
Gaussian on a linear continuum. The observed dispersion
$\sigma_{\rm obs}$ is corrected for the instrumental line-spread
function, $\sigma_{\rm inst}=\lambda_{\rm obs}/(R\,2.355)$ with $R=1000$
($\sigma_{\rm inst}\approx127\ \kms$), in quadrature:
$\sigma_{\rm int}^2=\sigma_{\rm obs}^2-\sigma_{\rm inst}^2$. We keep only
fits with residual ${\rm SNR}>5$ and intrinsic dispersion
$\sigma_{\rm int}<250\ \kms$, the latter cut removing broad-line AGN.
The circular velocity is taken to be
\begin{equation}
  \label{eq:vobs}
  \Vobs=\sqrt{2}\,\sigma_{\rm int},
\end{equation}
the nominal isotropic virialization factor for dispersion-dominated
systems. The sensitivity of our results to this choice is examined in
\S\ref{sec:discussion}.

\subsection{Baryonic masses}
\label{sec:masses}

The JADES DR4 external catalogue does not provide SED-fit stellar
masses, so we estimate $\Mbar$ from the \Ha{} luminosity through a
star-formation-rate (SFR) ladder. We convert the continuum-subtracted
\Ha{} flux (correcting for [N\,\textsc{ii}] where available) into a
luminosity using the luminosity distance, and adopt the
\citet{Kennicutt1998} calibration with a Chabrier-IMF factor of $0.63$,
$\mathrm{SFR}=0.63\,L_{\Ha}/1.26\times10^{41}\ \Msun\,{\rm yr^{-1}}$. We
then invert the star-forming main sequence of \citet{Popesso2023},
\begin{equation}
  \label{eq:ms}
  \log\mathrm{SFR}=0.84\log M_\star-8.2+1.1\log(1+z),
\end{equation}
to obtain $M_\star$, and add the gas reservoir with an evolving gas
fraction following \citet{Tacconi2020},
$\Mbar=M_\star/(1-f_{\rm gas})$. We retain galaxies with
$7<\log M_\star<12$. As an independent mass estimator we also compute
$\Mbar$ from $M_{\rm UV}$ via the \citet{Song2016}
$M_\star$--$M_{\rm UV}$ relation; the comparison of the two ladders is
central to the interpretation (\S\ref{sec:discussion}).

%% ============================================================
\section{Results}
\label{sec:results}

\subsection{Recovery of the $\beta(\Phi)$ shape}

For each galaxy we compute $\Vpred$ from equation (\ref{eq:btfr}) and the
deviation $\Phi$ from equation (\ref{eq:Phi}). Of the $59$ galaxies,
$24$ fall in the Taylor regime $|\Phi|<1$ where the linearized flow
applies; the remainder have large $|\Phi|$ driven by the steep
$V^4$ dependence and the scatter of the SFR-based masses, and are
excluded from the quantitative test. Figure~\ref{fig:beta} shows the
$24$ Taylor-regime galaxies in the $(\Phi,\beta)$ plane together with
the exact $\beta(\Phi)$ of equation (\ref{eq:beta}) and its quadratic
Taylor approximation. The data track the analytic curve: the
\emph{functional form} of the \LEG{} flow is recovered.

\begin{figure}[!t]
  \centering
  \includegraphics[width=\columnwidth]{fig2_beta}
  \caption{The $\beta$-function test. Points are the $24$ Taylor-regime
    JADES DR4 galaxies, colour-coded by redshift; the solid curve is the
    exact $\beta(\Phi)$ of equation (\ref{eq:beta}) and the dashed curve
    its quadratic approximation $-\Phi-\tfrac12\Phi^2$. The analytic
    shape predicted by \LEG{} is reproduced by the data.}
  \label{fig:beta}
\end{figure}

\subsection{A systematic offset towards $\Phi<0$}

While the shape is recovered, the population is not centred on the fixed
point. Figure~\ref{fig:vobs} compares $\Vobs$ with $\Vpred$: the
galaxies lie systematically \emph{below} the one-to-one line, with a
median ratio $\Vobs/\Vpred=0.77$. Equivalently, the Taylor-regime sample
has mean $\langle\Phi\rangle=-0.35$ and median $\Phi=-0.65$, and a
one-sample $t$-test rejects $\langle\Phi\rangle=0$ at $p=0.014$. The
measured intrinsic scatter is $\sigma_\Phi(0)=0.10$ dex
(Figure~\ref{fig:phiz}), i.e.\ a factor $\sim2.5$ larger than the
$<0.04$ dex predicted by equation (\ref{eq:sigmaPhi}). Thus the
\emph{magnitude} test is not passed: the high-$z$ BTFR residuals are
both offset and broader than the \LEG{} forecast.

\begin{figure}[!t]
  \centering
  \includegraphics[width=\columnwidth]{fig3_vobs_vpred}
  \caption{Observed circular velocity $\Vobs=\sqrt{2}\,\sigma_{\Ha}$
    versus the \LEG{} prediction $\Vpred=(G\,\Mbar\,\gcrit(z))^{1/4}$.
    The dashed line is one-to-one; the red line is the median scaling
    $\Vobs=0.77\,\Vpred$. Points are colour-coded by redshift.}
  \label{fig:vobs}
\end{figure}

\begin{figure}[!t]
  \centering
  \includegraphics[width=\columnwidth]{fig4_phi_z}
  \caption{BTFR deviation $\Phi$ versus redshift for the Taylor-regime
    sample. Shaded bands show the \LEG{} prediction
    $\pm\sigma_\Phi(0)\,H(z)/H_0$ for $\sigma_\Phi(0)=0.04$ dex (the
    \S10.8 target) and $0.10$ dex (measured). The data are offset to
    $\Phi<0$ and scatter beyond the $0.04$ dex band.}
  \label{fig:phiz}
\end{figure}

%% ============================================================
\section{Is the offset astrophysical? The SMITH-QA test}
\label{sec:smith}

A natural worry is that the offset is an artefact of equation
(\ref{eq:vobs}): if the \Ha{} dispersions of these $z=3$--$6$ starbursts
are inflated by turbulence, outflows or unresolved multiple components,
then $\sqrt{2}\,\sigma_{\Ha}$ would \emph{overestimate} the circular
velocity and bias $\Phi$. We test this with a model-independent
line-profile quality analysis, which we refer to as SMITH-QA. For each
of the $59$ galaxies we measure, directly from the residual of the
Gaussian fit, the skewness, excess kurtosis, reduced $\chi^2$, the
fraction of flux in the line wings, and the number of peaks. No object
is removed at this stage; the diagnostics only \emph{flag} non-Gaussian
profiles.

We then re-run the test after excluding the clearly non-Gaussian
profiles, defined by the conservative criterion
$\chi^2_{\rm red}>4$ \emph{or} $|\mathrm{skew}|>1$, which removes $13$ of
the $59$ galaxies. The result is summarized in
Table~\ref{tab:smith}: the median $\Vobs/\Vpred$ and median $\Phi$ are
\emph{unchanged} to three significant figures, and the scatter
$\sigma_\Phi(0)$ does not improve (it marginally worsens). The
systematic offset therefore does \emph{not} live in the shape of the
emission lines; cleaning the sample of turbulent or asymmetric profiles
leaves it intact. We note that $41$ of $59$ profiles formally show more
than one peak, but at $R\sim1000$ and with \Ha{} sampled by only
$\sim5$--$10$ pixels these are noise bumps rather than resolved
kinematic components: the $1$D integrated spectra simply do not carry the
information needed to separate rotation from turbulence. This is
precisely the regime in which \citet{FigueroaTorres2026} anticipated
that a clean dynamical measurement would not be possible.

\begin{table}[!t]\centering
  \setlength{\tabnotewidth}{0.95\columnwidth}
  \tablecols{3}
  \setlength{\tabcolsep}{1.6\tabcolsep}
  \caption{Robustness of the offset to line-profile quality cuts
    (SMITH-QA)} \label{tab:smith}
  \begin{tabular}{lcc}
    \toprule
    Quantity & Full ($59$) & After cut ($46$) \\
    \midrule
    median $\Vobs/\Vpred$      & $0.766$  & $0.766$  \\
    median $\Phi$              & $-0.649$ & $-0.649$ \\
    $\sigma_\Phi(0)$ [dex]     & $0.099$  & $0.106$  \\
    $p(\langle\Phi\rangle=0)$  & $0.014$  & $0.034$  \\
    \bottomrule
    \tabnotetext{}{Cut: $\chi^2_{\rm red}>4$ \emph{or}
      $|\mathrm{skew}|>1$, removing $13$ non-Gaussian profiles. The
      central values are unchanged; the offset is not an artefact of
      line-profile contamination.}
  \end{tabular}
\end{table}

%% ============================================================
\section{Discussion}
\label{sec:discussion}

Having excluded line-profile contamination, three interpretations of the
$\Phi<0$ offset remain: (A) the virialization factor in equation
(\ref{eq:vobs}); (B) a bias in the baryonic masses; or (C) a genuine
$\sim30\%$ tension with \LEG. We address each in turn.

\textit{(A) Virialization factor.} Replacing $\sqrt2$ by any constant
$\alpha\in[0.8,1.73]$ rescales $\Vobs$ but cannot remove the offset for a
plausible value: bringing the median to $\Phi=0$ would require
$\alpha=\sqrt2/0.77\approx1.8$, larger than the isotropic upper bound,
and a $t$-test still rejects $\langle\Phi\rangle=0$ for every $\alpha$ in
the range. The offset is not a virialization effect.

\textit{(B) Mass calibration.} This is the dominant systematic. To move
the median to $\Phi=0$ the baryonic masses would have to be reduced by
$-\log_{10}(1+\mathrm{median}\,\Phi)=+0.45$ dex (a factor $2.8$).
Crucially, this is \emph{smaller} than the dispersion between independent
mass estimators: for the same galaxies the \Ha$\to$SFR$\to$MS masses and
the $M_{\rm UV}\to M_\star$ masses differ by a median of $0.87$ dex
(Figure~\ref{fig:mass}). The signal we are trying to measure ($0.45$
dex) is buried inside the systematic floor of high-$z$ stellar masses
($0.87$ dex). In other words, the test is \emph{calibration-limited}:
with scaling-relation masses one cannot decide whether \LEG{} is right.

\textit{(C) Genuine tension.} We cannot exclude option (C), but neither
can the data demand it, because (B) alone spans the offset. A definitive
statement requires (i) SED-fit masses from the NIRCam photometry and
(ii) spatially resolved kinematics that measure $V_{\rm circ}$ directly
rather than through $\sqrt2\,\sigma$.

\begin{figure}[!t]
  \centering
  \includegraphics[width=\columnwidth]{fig5_mass}
  \caption{Inter-method dispersion of baryonic masses. The
    \Ha$\to$SFR$\to$main-sequence masses (used in the test) versus the
    $M_{\rm UV}\to M_\star$ masses for the same galaxies. The median
    offset is $0.87$ dex (grey band), larger than the $0.45$ dex mass
    reduction that would null the $\Phi<0$ offset. The systematic floor
    of high-$z$ masses exceeds the signal.}
  \label{fig:mass}
\end{figure}

This outcome is consistent with the framework's own forecast.
\citet{FigueroaTorres2026} note in \S10.8 that for samples at these
redshifts the relevant correction is below the current measurement floor
and becomes accessible only with extremely large telescopes; our
analysis quantifies that floor and identifies its origin as the mass
calibration rather than the kinematics. We stress what the test
\emph{does} establish: the analytic shape of $\beta(\Phi)$
(equation~\ref{eq:beta}) is reproduced by the Taylor-regime data
(Figure~\ref{fig:beta}), and the qualitative $\sigma_\Phi(z)\propto
H(z)/H_0$ scaling is not contradicted. It is the absolute normalization
of the scatter, and the sign of the mean deviation, that lie beyond the
reach of $R\sim1000$, scaling-relation-mass data.

%% ============================================================
\section{Conclusions}
\label{sec:conclusions}

We have carried out the first quantitative confrontation of the
Lambda Emergent Gravity $\beta(\Phi)$ flow with high-redshift
spectroscopy, using $59$ JADES DR4 galaxies at $z=3$--$6$. Our findings
are:

\begin{compactenum}[(1)]
\item The analytic shape of the \LEG{} $\beta$-function (equation
  \ref{eq:beta}) is recovered by the $24$ galaxies in the Taylor regime
  $|\Phi|<1$ (Figure~\ref{fig:beta}).
\item The population is offset from the fixed point, with median
  $\Vobs/\Vpred=0.77$, median $\Phi=-0.65$ ($p=0.014$), and a scatter
  $\sigma_\Phi(0)=0.10$ dex, $\sim2.5\times$ the predicted $<0.04$ dex.
\item A model-independent line-profile quality test (SMITH-QA) shows the
  offset is robust to the removal of non-Gaussian profiles
  (Table~\ref{tab:smith}): it is \emph{not} caused by turbulence or
  outflows in $\sigma_{\Ha}$.
\item The offset corresponds to a $+0.45$ dex mass overestimate, which
  lies within the $0.87$ dex inter-method scatter of high-$z$ masses
  (Figure~\ref{fig:mass}). The test is therefore
  \emph{calibration-limited}, not theory-limited, and does not by itself
  confirm or refute \LEG.
\end{compactenum}

A decisive test will require SED-fit baryonic masses and spatially
resolved kinematics at $z=4$--$6$, as will become available with
ELT/HARMONI and JWST/NIRSpec IFU programmes. The methodology developed
here---the $\Phi$/$\beta$ estimator and the SMITH-QA robustness
test---provides a template for that future measurement.

\acknowledgments
This work made use of public data from the JWST Advanced Deep
Extragalactic Survey (JADES) and the Mikulski Archive for Space
Telescopes (MAST). We thank the JADES team for making the DR4 catalogues
and spectra available.

\bibliographystyle{aasjournal-rm}
\bibliography{lambda_eg_jades}

\end{document}
